%%%%%%%%%%%%%%%%%%%%%%%%%%%%%%%%%%%%%%%%%%%%%%%%%%%%%%%%%%%%%%%%%%%%%%%%%%%%%%%%%%%%%%%%%%%%%%%%%%%%%%%%%%%%%%%%%%%%%%%%%%%%%%%%%%%%%
\subsection{Доказательство Утверждения 1}
%\textbf{Доказательство Утверждения 1:}\nopagebreak[4]

\begin{enumerate}

\item Рассмотрим произвольную функцию $g^\rho(\lambda,\Lambda,x,y)$ и подставим в нее представление (\ref{predst}). Раскладывая в многомерный ряд Тейлора получаем:

$$g^\rho(\lambda,\Lambda,x,y) = g^\rho(\lambda_0,\Lambda_0,x_0,y_0) + \sum_{j=1}^{+\infty} \frac{\mathcal{T}^j g^\rho}{j!},$$
где $\mathcal{T}$ - оператор вида:

\begin{multline*}
\mathcal{T} = 
\left(\sum_{k=1}^{+\infty} \lambda_k \varepsilon^k \right) \frac{\partial}{\partial \lambda} \bigg|_{\lambda_0,\Lambda_0,x_0,y_0}+
\left(\sum_{k=1}^{+\infty} \Lambda_k \varepsilon^k \right) \frac{\partial}{\partial \Lambda} \bigg|_{\lambda_0,\Lambda_0,x_0,y_0}+\\
+\left(\sum_{k=1}^{+\infty} x_k \varepsilon^k \right) \frac{\partial}{\partial x} \bigg|_{\lambda_0,\Lambda_0,x_0,y_0}
+\left(\sum_{k=1}^{+\infty} y_k \varepsilon^k \right) \frac{\partial}{\partial y} \bigg|_{\lambda_0,\Lambda_0,x_0,y_0}.
\end{multline*}

Вводя оператор
$$\mathfrak{D}_k = \left( \begin{pmatrix} \lambda_k\\ \Lambda_k \\ x_k \\ y_k \end{pmatrix},\begin{pmatrix} \frac{\partial}{\partial \lambda} |_{\lambda_0,\Lambda_0,x_0,y_0}\\ \frac{\partial}{\partial \Lambda} |_{\lambda_0,\Lambda_0,x_0,y_0} \\ \frac{\partial}{\partial x} |_{\lambda_0,\Lambda_0,x_0,y_0} \\ \frac{\partial}{\partial y} |_{\lambda_0,\Lambda_0,x_0,y_0} \end{pmatrix} \right),$$
получаем:
$$\mathcal{T} = \sum_{k=1}^{+\infty} \varepsilon^k \mathfrak{D}_k.$$

%$(\cdot,\cdot)$ - стандартное скалярное произведение в $\mathbb{R}^4$.

Рассмотрим выражение для $\mathcal{T}^j$ в котором, явно взяв степень, получаем:

$$\mathcal{T}^j = \left( \sum_{k=1}^{+\infty} \varepsilon^k \mathfrak{D}_k \right)^j = \sum_{k=j}^{+\infty} \varepsilon^k \mathcal{L}_k^{(j)},$$
где:
$$\mathcal{L}_k^{(j)} = \sum_{k_1+...+k_j = k} \mathfrak{D}_{k_1} \cdot ... \cdot \mathfrak{D}_{k_j}.$$

Подставляя этот результат в ряд Тейлора и группируя слагаемые по степеням $\varepsilon$ получаем:

$$g^\rho(\lambda,\Lambda,x,y) = g^\rho(\lambda_0,\Lambda_0,x_0,y_0) + \sum_{k=1}^{+\infty} \varepsilon^k \left( \sum_{j=1}^k \frac{\mathcal{L}_k^{(j)} g^\rho}{j!} \right).$$


Разобьем внутреннюю сумму на 2 части, учитывая $\mathcal{L}_k^{(1)} = \mathfrak{D}_k$:

\begin{equation*}
\begin{aligned}
g^\rho(\lambda,\Lambda,x,y) 
    &= g^\rho(\lambda_0,\Lambda_0,x_0,y_0) + \varepsilon \mathfrak{D}_1 g^\rho + \sum_{k=2}^{+\infty} \varepsilon^k \left( \mathcal{L}_k^{(1)}g^\rho + \sum_{j=2}^k \frac{\mathcal{L}_k^{(j)} g^\rho}{j!} \right) = \\ 
    &= g^\rho(\lambda_0,\Lambda_0,x_0,y_0) + \varepsilon \mathfrak{D}_1 g^\rho + \sum_{k=2}^{+\infty} \varepsilon^k \left( \mathfrak{D}_k g^\rho + \underbrace{\sum_{j=2}^k \frac{\mathcal{L}_k^{(j)} g^\rho}{j!}}_{G_k^\rho} \right) \\ 
    %&= g(\lambda_0,\Lambda_0,x_0,y_0) + \underbrace{\sum_{k=1}^{\infty} \varepsilon^k \mathfrak{D}_k g}_{I} 
    %+ \underbrace{\sum_{k=2}^{\infty} \varepsilon^k \left( \sum_{j=2}^{k} \frac{C_k^{(j)} g }{j!} \right)}_{II = G_k^\rho}.
\end{aligned}
\end{equation*}

\begin{itemize}
\item $\mathfrak{D}_k g^\rho$ - содержит $(\lambda_k, \Lambda_k, x_k, y_k)$ и уйдет в оператор однородного уравнения;
\item $G_k^\rho$ - содержит $(\lambda_j, \Lambda_j, x_j, y_j), \quad j \le k-1$ и входит в неднородность в уравнениях (\ref{xy_eq}).
\end{itemize}

Введенные таким образом $G_k^\rho$ определены только для $k \ge 2$, для $k=1$ по определению положим $G_1^\rho = 0$. В итоге получаем:

\begin{equation}
G_k^\rho = \sum_{j=2}^{k} \frac{\mathcal{L}_k^{(j)} g^\rho }{j!}, \quad k \ge 2,
\label{G_def}
\end{equation}
$$G_1^\rho \equiv 0,$$
$$g^\rho(\lambda,\Lambda,x,y) = g^\rho(\lambda_0,\Lambda_0,x_0,y_0) + \sum_{k=1}^{\infty} \varepsilon^k \mathfrak{D}_k g^\rho + \sum_{k=2}^{+\infty} \varepsilon^k G_k^\rho.$$

\item Так как $g^\rho, \rho \in \{\lambda, \Lambda, x, y\}$ и все ее производные ограничены, а сумма (\ref{G_def}) начинается с $j=2$, то $G_{k+1}^\rho$ будет содержать слагаемые из 2 и более множителей (порядка $\le k$) вида $\mathfrak{D}_i \mathfrak{D}_j g^\rho, i+j \leq k$. Учитывая что каждое решение порядка $\le k$ принадлежит классу $\mathcal{F}_s$, то их произведения и, как следствие, $G_{k+1}^\rho$ будут принадлежать классу $\mathcal{F}_{2s}$.

\item Так как $g^\lambda = \alpha \Lambda$ не содержит нелинейные члены, то любые производные второго порядка и старше будут равны 0. В следствие этого $G_k^\lambda = 0 \quad \forall k \ge 1$. $\blacksquare$
\end{enumerate}
%%%%%%%%%%%%%%%%%%%%%%%%%%%%%%%%%%%%%%%%%%%%%%%%%%%%%%%%%%%%%%%%%%%%%%%%%%%%%%%%%%%%%%%%%%%%%%%%%%%%%%%%%%%%%%%%%%%%%%%%%%%%%%%%%%%%%






\subsection{Об однородности и неоднородности уравнений первого приближения}
%%%%%%%%%%%%%%%%%%%%%%%%%%%%%
%\textbf{Доказательство:}\nopagebreak[4]
\begin{enumerate}

\item Рассмотрим абстрактную 4 мерную быстро-медленную систему:

\begin{equation*}
    \begin{cases}
        \dot \Lambda(t)=g^\Lambda(\Lambda,\lambda,x,y), \\
        \dot \lambda(t)=g^\lambda(\Lambda,\lambda,x,y), \\
        \dot x(t)=\varepsilon g^x(\Lambda,\lambda,x,y), \\
        \dot y(t)=\varepsilon g^y(\Lambda,\lambda,x,y), \\
    \end{cases}
\end{equation*}

Переходя к медленному времени $\tau = \varepsilon t$ получаем:
\begin{equation*}
    \begin{cases}
        \varepsilon \Lambda'(\tau)=g^\Lambda(\Lambda,\lambda,x,y), \\
        \varepsilon \lambda'(\tau)=g^\lambda(\Lambda,\lambda,x,y), \\
        x'(\tau)=g^x(\Lambda,\lambda,x,y), \\
        y'(\tau)=g^y(\Lambda,\lambda,x,y), \\
    \end{cases}
\end{equation*}

Обозначим $\v X = (\Lambda,\lambda,x,y)$ и будем искать формальное решение в виде:
$\v X = \sum_{k=0}^\infty \varepsilon^k \v X_k.$
Подставляя это представление в систему и приравнивая члены при одинаковых степенях $\varepsilon$ получаем для младших порядков:

\begin{equation}
    \begin{cases}
        0=g^\Lambda(\v X_0), \\
        0=g^\lambda(\v X_0), \\
        x_0'(\tau)=g^x(\v X_0), \\
        y_0'(\tau)=g^y(\v X_0), \\
    \end{cases}
\end{equation}

\begin{equation}
    \begin{cases}
        \Lambda_0'(\tau)=(\nabla g^\Lambda(\v X_0),\v X_1), \\
        \lambda_0'(\tau)=(\nabla g^\lambda(\v X_0),\v X_1), \\
        x_1'(\tau)=(\nabla g^x(\v X_0),\v X_1), \\
        y_1'(\tau)=(\nabla g^y(\v X_0),\v X_1), \\
    \end{cases}
    \label{1st_order}
\end{equation}

Заметим, что первые два уравнения в (\ref{1st_order}) являются линейной неоднородной системой алгебраических уравнений по переменным $(\Lambda_1, \lambda_1)$, представимой в виде:

\begin{align*}
\begin{pmatrix}
(\nabla g^\Lambda)_\Lambda & (\nabla g^\Lambda)_\lambda \\
(\nabla g^\lambda)_\Lambda & (\nabla g^\lambda)_\lambda
\end{pmatrix}
\begin{pmatrix} \Lambda_1 \\ \lambda_1 \end{pmatrix}
= A \begin{pmatrix} \Lambda_1 \\ \lambda_1 \end{pmatrix} =
\begin{pmatrix} \Lambda_0' - (\nabla g^\Lambda)_x x_1 - (\nabla g^\Lambda)_y y_1 \\ \lambda_0' - (\nabla g^\lambda)_x x_1 - (\nabla g^\lambda)_y y_1 \end{pmatrix},
\end{align*}

\begin{equation*}
A = \begin{pmatrix}
(\nabla g^\Lambda)_\Lambda & (\nabla g^\Lambda)_\lambda \\
(\nabla g^\lambda)_\Lambda & (\nabla g^\lambda)_\lambda
\end{pmatrix}.
\end{equation*}

Тогда решение этой системы имеет вид:
\begin{align*}
    \begin{pmatrix} \Lambda_1 \\ \lambda_1 \end{pmatrix} = 
    M \begin{pmatrix} x_1 \\ y_1 \end{pmatrix} + \v C,
\end{align*}

\begin{align*}
    M = \frac{1}{\det A} \begin{pmatrix} 
        -(\nabla g^\lambda)_\lambda (\nabla g^\Lambda)_x + (\nabla g^\Lambda)_\lambda (\nabla g^\lambda)_x & -(\nabla g^\lambda)_\lambda (\nabla g^\Lambda)_y + (\nabla g^\Lambda)_\lambda (\nabla g^\lambda)_y \\ 
         (\nabla g^\lambda)_\Lambda (\nabla g^\Lambda)_x - (\nabla g^\Lambda)_\Lambda (\nabla g^\lambda)_x &  (\nabla g^\lambda)_\Lambda (\nabla g^\Lambda)_y - (\nabla g^\Lambda)_\Lambda (\nabla g^\lambda)_y
    \end{pmatrix},
\end{align*}

$$\v C = \frac{1}{\det A} \begin{pmatrix} (\nabla g^\lambda)_\lambda \Lambda_0' - (\nabla g^\Lambda)_\lambda \lambda_0' \\ -(\nabla g^\lambda)_\Lambda \Lambda_0' + (\nabla g^\Lambda)_\Lambda \lambda_0' \end{pmatrix}.$$

Подставляя решение в последние два уравнения в (\ref{1st_order}) и группируя слагаемые получаем:

\begin{equation}
    \begin{cases}
        x_1'(\tau) = 
            \Big( 
                (\nabla g^x)_x + \big( (\nabla g^x)_\Lambda M_{11} + (\nabla g^x)_\lambda M_{21} \big)
            \Big)x_1 +\\
            \quad \quad \quad + \Big(
                (\nabla g^x)_y + \big( (\nabla g^x)_\Lambda M_{12} + (\nabla g^x)_\lambda M_{22} \big)
            \Big)y_1 +\\
            \quad \quad \quad + \left( \begin{pmatrix} (\nabla g^x)_\Lambda \\ (\nabla g^x)_\lambda \end{pmatrix},\v C \right), \\
            \\
        y_1'(\tau)=
            \Big(
                (\nabla g^y)_x + \big( (\nabla g^y)_\Lambda M_{11} + (\nabla g^y)_\lambda M_{21} \big)
            \Big)x_1 +\\
            \quad \quad \quad + \Big(
                (\nabla g^x)_y + \big( (\nabla g^y)_\Lambda M_{12} + (\nabla g^y)_\lambda M_{22} \big)
            \Big)y_1 +\\
            \quad \quad \quad + \left( \begin{pmatrix} (\nabla g^y)_\Lambda \\ (\nabla g^y)_\lambda \end{pmatrix},\v C \right). \\
    \end{cases}
    \label{x1y1_eq}
\end{equation}

Здесь $(\cdot,\cdot)$ -- стандартное скалярное произведение в $\mathbb{R}^n$.

Таким образом, в общем случае, уравнение для первого приближения является неоднородным, однако для рассматриваемой системы (\ref{fulltn}) неоднородность равна нулю. Действительно, заметим:

$$\Lambda_0 = 0 \quad \Rightarrow \quad \v C = \begin{pmatrix} \frac{\lambda_0'}{\alpha} \\ 0 \end{pmatrix}$$

Также заметим, что компоненты правой части $g^x, g^y$ не зависят от $\Lambda$, поэтому векторы, входящие в неоднородность в уравнениях (\ref{x1y1_eq}) имеют вид:

$$\begin{pmatrix} (\nabla g^x)_\Lambda \\ (\nabla g^x)_\lambda \end{pmatrix} = \begin{pmatrix} 0 \\ * \end{pmatrix},$$
$$\begin{pmatrix} (\nabla g^y)_\Lambda \\ (\nabla g^y)_\lambda \end{pmatrix} = \begin{pmatrix} 0 \\ * \end{pmatrix}.$$

Тогда неоднородность в системе (\ref{x1y1_eq}) вырождается:

$$\left( \begin{pmatrix} (\nabla g^x)_\Lambda \\ (\nabla g^x)_\lambda \end{pmatrix},\v C \right) = \left( \begin{pmatrix} (\nabla g^y)_\Lambda \\ (\nabla g^y)_\lambda \end{pmatrix},\v C \right) = 0.$$

Таким образом, в системе (\ref{fulltn}) дифференциальное уравнение на $(x_1,y_1)$ является однородным.

\item 

Рассмотрим уравнения медленной системы (\ref{slow_sys}). Для простоты приведем выкладки только для компоненты $x_0$.

Продифференцируем уравнение на $x_0$ по времени, при этом будем считать, что $\Lambda_0, \lambda_0$ зависят от $x_0,y_0$:
\begin{multline}
x_0'' = (\nabla g^x)_x x_0' + (\nabla g^x)_y y_0' + (\nabla g^x)_\Lambda \Lambda_0' + (\nabla g^x)_\lambda \lambda_0' = \\
= 
\left( (\nabla g^x)_x + (\nabla g^x)_\Lambda\frac{\partial \Lambda_0}{\partial x} + (\nabla g^x)_\lambda\frac{\partial \lambda_0}{\partial x} \right) x_0' +\\
+\left( (\nabla g^x)_y + (\nabla g^x)_\Lambda\frac{\partial \Lambda_0}{\partial y} + (\nabla g^x)_\lambda\frac{\partial \lambda_0}{\partial y} \right) y_0'.
\label{xdd_eq}
\end{multline}

Заметим, что $\Lambda_0(x_0,y_0), \lambda_0(x_0,y_0)$ задаются неявно через уравнения:
\begin{equation}
    \begin{cases}
        0=g^\Lambda(\v X_0), \\
        0=g^\lambda(\v X_0). \\
    \end{cases}
\end{equation}

Дифференцируем оба уравнения по \(x\):
\[
\begin{cases}
(\nabla g^\Lambda)_x + (\nabla g^\Lambda)_y \underbrace{y_x}_0 + (\nabla g^\Lambda)_\Lambda \cdot \Lambda_x + (\nabla g^\Lambda)_\lambda \cdot \lambda_x = 0, \\
(\nabla g^\lambda)_x + (\nabla g^\lambda)_y \underbrace{y_x}_0 + (\nabla g^\lambda)_\Lambda \cdot \Lambda_x + (\nabla g^\lambda)_\lambda \cdot \lambda_x = 0.
\end{cases}
\]

Полученное выражение можно переписать в матричной форме:
\[
\begin{pmatrix}
(\nabla g^\Lambda)_\Lambda & (\nabla g^\Lambda)_\lambda \\
(\nabla g^\lambda)_\Lambda & (\nabla g^\lambda)_\lambda
\end{pmatrix}
\begin{pmatrix}
\Lambda_x \\
\lambda_x
\end{pmatrix}
= -
\begin{pmatrix}
(\nabla g^\Lambda)_x \\
(\nabla g^\lambda)_x
\end{pmatrix}.
\]

Решая это уравнение относительно $(\Lambda_x,\lambda_x)$ получаем:
$$
\Lambda_x = \frac{(\nabla g^\Lambda)_\lambda (\nabla g^\lambda)_x - (\nabla g^\Lambda)_x (\nabla g^\lambda)_\lambda}{(\nabla g^\Lambda)_\Lambda (\nabla g^\lambda)_\lambda - (\nabla g^\Lambda)_\lambda (\nabla g^\lambda)_\Lambda},
$$
$$
\lambda_x = \frac{(\nabla g^\Lambda)_x (\nabla g^\lambda)_\Lambda - (\nabla g^\Lambda)_\Lambda (\nabla g^\lambda)_x}{(\nabla g^\Lambda)_\Lambda (\nabla g^\lambda)_\lambda - (\nabla g^\Lambda)_\lambda (\nabla g^\lambda)_\Lambda}.
$$

Заметим, что данные частные производные в точности совпадают с коэффициентами матрицы $M$:
$$
\frac{\partial \Lambda_0}{\partial x} = M_{11},
$$
$$
\frac{\partial \lambda_0}{\partial x} = M_{21},
$$
Аналогично можно показать следующее:
$$
\frac{\partial \Lambda_0}{\partial y} = M_{12},
$$
$$
\frac{\partial \lambda_0}{\partial y} = M_{22}.
$$

Подставляя данный результат в (\ref{xdd_eq}) получаем:
\begin{multline}
(x_0')' = (\nabla g^x)_x x_0' + (\nabla g^x)_y y_0' + (\nabla g^x)_\Lambda \Lambda_0' + (\nabla g^x)_\lambda \lambda_0' = \\
= 
\left( (\nabla g^x)_x + (\nabla g^x)_\Lambda M_{11} + (\nabla g^x)_\lambda M_{21} \right) x_0' +\\
+\left( (\nabla g^x)_y + (\nabla g^x)_\Lambda M_{12} + (\nabla g^x)_\lambda M_{22} \right) y_0'.
\label{xdd_eq2}
\end{multline}

Учитывая, что в системе (\ref{fulltn}) неоднородность в уравнениях (\ref{x1y1_eq}) вырождается, вид уравнений (\ref{xdd_eq2}) и (\ref{x1y1_eq}) полностью совпадает.

Заметим, что асимптотика сепаратрисы медленной системы при $\tau \to \pm \infty$ имеет вид: 
\begin{equation*}
    \begin{dcases}
        x_0(\tau) = \frac{\text{const}}{\cosh(f_\pm \cdot \tau)} + \mathcal{O}(e^{-2 f_\pm |\tau|}), \quad \tau \to \pm \infty,\\
        y_0(\tau) = \frac{\text{const} \cdot \sinh(f_\pm \cdot \tau)}{\cosh^2(f_\pm \cdot \tau)} + \mathcal{O}(e^{-2 f_\pm |\tau|}), \quad \tau \to \pm \infty,
    \end{dcases}
\end{equation*}
где $f_\pm = \lim_{\tau \to \pm \infty} \frac{d f^\pm(\tau)}{d\tau} = \frac{J^{+} \chi^{+}}{\sigma^{+}}  \lim_{\tau \to \pm \infty} \frac{P^{+}(f^{+}(\tau))}{S^{+}(f^{+}(\tau))}$. Так как $P^\pm$ и $S^\pm$ являются полиномами второй степени от $\cosh f^\pm$, то пределы существуют и  равны:
 
$$f_+ = f_- = \frac{J^{+} \chi^{+}}{\sigma^{+}} \cdot \frac{1}{U_0} >0$$

Тогда получаем:
$$x_0,y_0 \in \mathcal{F}_{s_0}, \quad s_0 = f_+ = f_- > 0$$

Решение для первого приближения имеет аналогичное поведение на обеих бесконечностях:
\begin{equation*}
    \begin{cases}
        x_1(\pm \infty) = 0, \\
        
        y_1(\pm \infty) = 0. \\
    \end{cases}
\end{equation*}
$$x_1, y_1 \in \mathcal{F}_{s_0}$$

Отсюда получаем, что $(x_1,y_1)$ удовлетворяет граничным условиям и для $W^s(0)$, и для $W^u(0)$, поэтому индексы $s,u$ будем для него опускать. 

$\blacksquare$

\end{enumerate}%




%%%%%%%%%%%%%%%%%%%%%%%%%%%%%%%%%%%%%%%%%%%%%%%%%%%%%%%%%%%%%%%%%%%%
\begin{utv}

Для $G_k^\rho$ верна рекуррентная формула:

\begin{equation*}
G_{k+1}^\rho = \frac12 \sum_{m=1}^k \mathfrak{D}_m \mathfrak{D}_{k+1-m} g^\rho + \sum_{m=1}^k \mathfrak{D}_m \hat G_{k+1-m}^\rho,
\end{equation*}
где $\hat G_{k+1}^\rho$ отличается от $G_{k+1}^\rho$ нормировкой:
$$G_k^\rho = \sum_{j=2}^{k} \frac{\mathcal{L}_k^{(j)} g^\rho }{j!},$$
$$\hat G_k^\rho = \sum_{j=2}^{k} \frac{\mathcal{L}_k^{(j)} g^\rho }{(j+1)!},$$
$$\rho \in {\lambda, \Lambda, x, y}.$$
\end{utv}
%%%%%%%%%%%%%%%%%%%%%%%%%%%%%%%%%%%%%%%%%%%%%%%%%%%%%%%%%%%%%%%%%%%%%%%%%%%%%%%%%%%%%%%%%%%%%%%%%%%%%%%%%%%%%%%%%%%%%%%%%%%%%%%%%%%%%
\textbf{Доказательство:}\nopagebreak[4]
$$\mathcal{L}_{k+1}^{(j)}g^\rho = \sum_{m=1}^{k+2-j} \mathfrak{D}_m \mathcal{L}_{k+1-m}^{(j-1)} g^\rho$$

$$G_{k+1}^\rho = \sum_{j=2}^{k+1} \frac{\mathcal{L}_{k+1}^{(j)} g^\rho }{j!} = \sum_{j=2}^{k+1} \sum_{m=1}^{k+2-j} \frac{\mathcal{L}_{k+1-m}^{(j-1)} g^\rho }{j!} = \sum_{m=1}^k \mathfrak{D}_m \sum_{j=2}^{k+2-m} \frac{\mathcal{L}_{k+1-m}^{(j-1)} g^\rho }{j!} =$$
$$= \sum_{m=1}^k \mathfrak{D}_m \sum_{l=1}^{k+1-m} \frac{\mathcal{L}_{k+1-m}^{(l)} g^\rho }{(l+1)!} = \sum_{m=1}^k \mathfrak{D}_m \left( \sum_{l=2}^{k+1-m} \frac{\mathcal{L}_{k+1-m}^{(l)} g^\rho }{(l+1)!} + \frac12 \mathcal{L}_{k+1-m}^{(1)}g^\rho \right) =$$
$$= \sum_{m=1}^k \mathfrak{D}_m \left( \hat G_{k+1-m}^\rho + \frac12 \mathfrak{D}_{k+1-m}g^\rho \right) \blacksquare$$



%%%%%%%%%%%%%%%%%%%%%%%%%%%%%%%%%%%%%%%%%%%%%%%%%%%%%%%%%%%%%%%%%%%%%%%%%%%%
\begin{utv}
\begin{enumerate}

\item 
Если $\lambda_j,\Lambda_j,x_j,y_j, \quad j \le k-1$ как функции $\tau$ имеют четность, то \newline $G_k^x(\tau),\tilde G_k^x(\tau),G_k^\Lambda(\tau), G_k^y(\tau)$ также имеют четность. При этом $G_k^x(\tau),\tilde G_k^x(\tau),G_k^\Lambda(\tau)$ имеют одинаковую четность, а $G_k^y(\tau), \tilde G_k^y(\tau)$ имеют противоположную к $G_k^x(\tau)$ четность. 

\item В частности $G_2^x(\tau), \tilde G_2^x(\tau), G_2^\Lambda$ - нечетные, $G_2^y(\tau), \tilde G_2^y(\tau)$ - четные.
\end{enumerate}
\end{utv}
%%%%%%%%%%%%%%%%%%%%%%%%%%%%%%%%%%%%%%%%%%%%%%%%%%%%%%%%%%%%%%%%%%%%%%%%%%%%%%%%%%%%%%%%%%%%%%%%%%%%%%%%%%%%%%%%%%%%%%%%%%%%%%%%%%%%%
\textbf{Доказательство:}\nopagebreak[4]
\begin{enumerate}
\item Рассмотрим четность правой части и различных ее производных по координатам в подстановке нулевого приближения.
Т.к. $x_0(\tau)$ -- четная, $y_0(\tau), \lambda_0(\tau)$ -- нечетные получаем:
$$g^\Lambda|_{\lambda_0(\tau),\Lambda_0(\tau),x_0(\tau),y_0(\tau)} = U(x_0,y_0)\sin\lambda_0-V(x_0,y_0)\cos\lambda_0 \quad \text{ - нечетная},$$
$$g^x|_{\lambda_0(\tau),\Lambda_0(\tau),x_0(\tau),y_0(\tau)} = -\varepsilon \left( 2Fy_0-\frac{\partial U}{\partial y} \cos \lambda_0 - \frac{\partial V}{\partial y} \sin \lambda_0 \right) \quad \text{ - нечетная},$$
$$g^y|_{\lambda_0(\tau),\Lambda_0(\tau),x_0(\tau),y_0(\tau)} = \varepsilon \left( 2F(x_0+\hat x_0)+e_JG -\frac{\partial U}{\partial x} \cos \lambda_0 - \frac{\partial V}{\partial x} \sin \lambda_0 \right) \quad \text{ - четная}.$$

Заметим:
\begin{itemize}
    \item Каждая производная по чётной переменной ($x_0$) сохраняет чётность.
    \item Каждая производная по нечётной переменной ($y_0$ или $\lambda_0$) меняет чётность на противоположную.
\end{itemize}
Таким образом, обозначая четность как $+1$ для четных функций и  $-1$ -  для нечетных функций, получаем:
$$\frac{\partial^{p+q+m}g^\Lambda}{\partial x^p \partial y^q \partial \lambda^m}\bigg|_{\lambda_0(\tau),\Lambda_0(\tau),x_0(\tau),y_0(\tau)} \quad \Leftrightarrow \quad -(-1)^{q+m},$$
\begin{equation}
\frac{\partial^{p+q+m}g^x}{\partial x^p \partial y^q \partial \lambda^m}\bigg|_{\lambda_0(\tau),\Lambda_0(\tau),x_0(\tau),y_0(\tau)} \quad \Leftrightarrow \quad -(-1)^{q+m}, \label{chetnost}
\end{equation}
$$\frac{\partial^{p+q+m}g^y}{\partial x^p \partial y^q \partial \lambda^m}\bigg|_{\lambda_0(\tau),\Lambda_0(\tau),x_0(\tau),y_0(\tau)} \quad \Leftrightarrow \quad (-1)^{q+m}.$$

Заметим, что $G_k^\Lambda(\tau),G_k^x(\tau),G_k^y(\tau)$ имеют вид:
$$G_k^\Lambda(\tau) = \sum_{p_1,q_1,m_1,...} \frac{\partial^{p+q+m}g^\Lambda}{\partial x^p \partial y^q \partial \lambda^m}\bigg|_{\lambda_0(\tau),\Lambda_0(\tau),x_0(\tau),y_0(\tau)} (x_{p_1}y_{q_1}\lambda_{m_1}...),$$
$$G_k^x(\tau) = \sum_{p_1,q_1,m_1,...} \frac{\partial^{p+q+m}g^x}{\partial x^p \partial y^q \partial \lambda^m}\bigg|_{\lambda_0(\tau),\Lambda_0(\tau),x_0(\tau),y_0(\tau)} (x_{p_1}y_{q_1}\lambda_{m_1}...),$$
$$G_k^y(\tau) = \sum_{p_1,q_1,m_1,...} \frac{\partial^{p+q+m}g^y}{\partial x^p \partial y^q \partial \lambda^m}\bigg|_{\lambda_0(\tau),\Lambda_0(\tau),x_0(\tau),y_0(\tau)} (x_{p_1}y_{q_1}\lambda_{m_1}...),$$
и отличаются только производными, множители в скобках у них одинаковые.

Следовательно, если определена четность множителя в скобках (т.е. определена четность всех решений порядка $\le k-1$), то четность $G_k^\Lambda(\tau),G_k^x(\tau),G_k^y(\tau)$ отличается на четность производных.
Учитывая, что четность производных одинакового порядка для $g^\Lambda$ и $g^x$ совпадает, а для $g^y$ отличается, получаем необходимое утверждение.

Четность $\tilde G_k^x, \tilde G_k^y$ наследуется от $G_k^\Lambda(\tau),G_k^x(\tau),G_k^y(\tau)$.

\item Заметим, что $x_0(\tau), y_1(\tau), \tilde x_1, \lambda_1$ -- четные, а $y_0(\tau), x_1(\tau), \tilde y_1$ -- нечетные.

\begin{equation}
G_2^\rho(\tau) = \sum_{p+q+m=2,\quad p,q,m \ge 0, \quad p,q,m \neq 2} \frac{\partial^{2}g^\rho}{\partial x^p \partial y^q \partial \lambda^m}\bigg|_{\lambda_0(\tau),\Lambda_0(\tau),x_0(\tau),y_0(\tau)} x_1^p y_1^q \lambda_1^m
\label{g_2}
\end{equation}

Кроме того, четность $x_1^p y_1^q \lambda_1^m$ равна $(-1)^p$. Тогда, учитывая (\ref{chetnost}), получаем что четность слагаемых в (\ref{g_2}) равна $-(-1)^{p+q+m}$ для $\rho \in \{ \Lambda, x \}$ и $(-1)^{p+q+m}$ для $\rho = y$. Но так как во всех слагаемых $p+q+m=2$, получаем:
$G_2^x(\tau), G_2^\Lambda$ - нечетные, $G_2^y(\tau)$ - четные.

Четность $\tilde G_2^x, \tilde G_2^y$ наследуется от $G_2^\Lambda(\tau),G_2^x(\tau),G_2^y(\tau)$.
$\blacksquare$
\end{enumerate}
%%%%%%%%%%%%%%%%%%%%%%%%%%%%%%%%%%%%%%%%%%%%%%%%%%%%%%%%%%%%%%%%%%%%%%%%%%%%%%%%%%%%%%%%%%%%%%%%%%%%%%%%%%%%%%%%%%%%%%%%%%%%%%%%%%%%%